\section{Introducción}

En un mercado turístico globalizado y de alta exigencia, la capacidad de las organizaciones para gestionar sus flujos de información de manera eficiente y segura se ha convertido en un factor determinante para su sostenibilidad y competitividad. En este contexto, Machu Picchu Exclusive Tours E.I.R.L., una agencia de viajes cusqueña fundada en 2015 y especializada en turismo receptivo de lujo y aventura, ha desarrollado su actividad basándose en el profundo conocimiento del patrimonio cultural peruano por parte de sus guías locales. Si bien esta especialización le ha permitido consolidar una reputación positiva en el servicio personalizado, la empresa operaba históricamente bajo un modelo de gestión con una limitada digitalización de sus procesos comerciales y administrativos.

Antes de la intervención tecnológica, la gestión operativa de la agencia se realizaba de forma manual, apoyándose predominantemente en herramientas como Excel y WhatsApp. Esta dependencia generaba una fragmentación de la data, donde la información de los pasajeros, las cotizaciones y los itinerarios ---que varían en duración de 2 a 20 días--- se encontraban dispersos, dificultando la trazabilidad y el seguimiento de las oportunidades comerciales. Un punto crítico identificado en el diagnóstico inicial fue el manejo de información sensible: la recepción de fotografías de pasaportes a través de dispositivos personales representaba un riesgo latente de ciberseguridad, obligando a la gerencia a depurar estos archivos periódicamente para proteger la privacidad del cliente, lo que a su vez impedía contar con una base de datos histórica estructurada.

Ante esta problemática y con el objetivo de escalar hacia los mercados internacionales norteamericano y europeo, se planteó la implementación del sistema ERP Odoo versión 19 en modalidad nube (SaaS). Esta solución integral fue seleccionada para centralizar la data operativa mediante el uso de funcionalidades estándar, asegurando una adopción ágil y segura sin necesidad de desarrollos a medida. La implementación abarcó módulos clave como el CRM Turístico para la gestión de relaciones y almacenamiento seguro de documentos, el módulo de Proyectos para el control visual de itinerarios mediante tableros Kanban, y la gestión de Ventas y Facturación para el registro manual de adelantos y el control riguroso de saldos pendientes antes de la compra de boletos no reembolsables.

El presente informe tiene como finalidad evaluar los resultados obtenidos tras la puesta en marcha del sistema Odoo. A través de este documento, se analiza el impacto generado en la operación diaria de la empresa, la optimización de sus procesos de captación y reserva, y el nivel de adopción tecnológica alcanzado por el equipo humano. El reporte evidencia cómo la transición de una gestión manual a un entorno digital centralizado ha permitido profesionalizar la atención al cliente, fortalecer la seguridad de la información y dotar a la gerencia de herramientas analíticas para una toma de decisiones basada en evidencia.
